\begin{solapartat}
Per tal que el fil es trobi suspés a l'aire, cal que hi hagi una força que compensi la força de gravetat. Per tant, el camp magnètic ha de fer una força sobre el fil cap amunt, és a dir, en la direcció positiva de l'eix $Z$. \punts{0,2} Partint de l'expressió $\vec{F} = I\vec{L}\wedge\vec{B}$ \punts{0,2} amb $\vec{F}$ \punts{0,2} apuntant cap a les $z$'s positives, i el vector $\vec{L}$ cap a les $x$'s positives, el camp $\vec{B}$ necessàriament ha d'apuntar en la direcció positiva de l'eix $Y$.

Pel que fa al seu mòdul, cal que compensi el de la força de gravetat, i per tant
\[ ILB = mg\ \text{\punts{0,2}} \]
d'on resulta
\[ B = \frac{mg}{IL} = \frac{0,1\times 9,8}{0,3\times 5} = 0,65\ \mathrm{T}\ \text{\punts{0,2}} \]
\end{solapartat}

\begin{solapartat}
Amb el fil rectilini fem una espira circular que situem al pla $XY$, i el camp magnètic també es troba situat al pla $XY$, de forma que el flux del camp és $\Phi = \vec{B}\cdot\vec{S} = 0$, en cada instant, ja que els vectors camp magnètic i superfície de l'espira són perpendiculars. \punts{1}
\end{solapartat}
